\textbf{Phần 1. Trắc nghiệm lựa chọn}
\setcounterpageref{ex}{0}

\begin{ex} %[1D3V11]
Tính kết quả của biểu thức: $10 - 3 - 2$.
\choice
{$4$}
{$6$}
{\True $5$}
{$3$}
\loigiai{
Thực hiện từ trái sang phải: $10 - 3 = 7$; sau đó $7 - 2 = 5$.
\textbf{Đáp số:} $5$.\par
\textbf{Đáp án của câu hỏi là C}
}
\end{ex}

\begin{ex} %[1D3V11]
Tính kết quả của biểu thức: $9 - 4 - 1$.
\choice
{$3$}
{\True $4$}
{$5$}
{$2$}
\loigiai{
Thực hiện từ trái sang phải: $9 - 4 = 5$; sau đó $5 - 1 = 4$.
\textbf{Đáp số:} $4$.\par
\textbf{Đáp án của câu hỏi là B}
}
\end{ex}

\begin{ex} %[1D3V11]
Điền số thích hợp vào ô trống: $8 - \dots = 3$.
\choice
{$4$}
{$6$}
{$3$}
{\True $5$}
\loigiai{
Muốn tìm số trừ, ta lấy số bị trừ trừ đi hiệu: $8 - 3 = 5$.
\textbf{Đáp số:} $5$.\par
\textbf{Đáp án của câu hỏi là D}
}
\end{ex}

\begin{ex} %[1D3V11]
Điền số thích hợp vào ô trống: $\dots - 4 = 5$.
\choice
{\True $9$}
{$8$}
{$10$}
{$7$}
\loigiai{
Muốn tìm số bị trừ, ta lấy hiệu cộng với số trừ: $5 + 4 = 9$.
\textbf{Đáp số:} $9$.\par
\textbf{Đáp án của câu hỏi là A}
}
\end{ex}

\begin{ex} %[1D3V11]
Điền số thích hợp: $10 - \dots = 4 + 2$.
\choice
{$3$}
{$5$}
{\True $4$}
{$2$}
\loigiai{
Vế phải: $4 + 2 = 6$.\nTa có: $10 - \dots = 6 \Rightarrow$ số cần điền là: $10 - 6 = 4$.
\textbf{Đáp số:} $4$.\par
\textbf{Đáp án của câu hỏi là C}
}
\end{ex}

\begin{ex} %[1D3V11]
Trên cành cây có 9 chú chim. Có 3 chú chim bay đi, lát sau lại có thêm 2 chú chim nữa bay đi. Hỏi trên cành cây còn lại bao nhiêu chú chim?
\choice
{$3$}
{\True $4$}
{$5$}
{$2$}
\loigiai{
Sau lần thứ nhất còn: $9 - 3 = 6$ (chú chim).\nSau lần thứ hai trên cành còn: $6 - 2 = 4$ (chú chim).
\textbf{Đáp số:} $4$.\par
\textbf{Đáp án của câu hỏi là B}
}
\end{ex}

\begin{ex} %[1D3V11]
Trong hộp có 10 chiếc kẹo. Nam ăn 2 chiếc kẹo, cho em 3 chiếc kẹo. Hỏi trong hộp còn lại bao nhiêu chiếc kẹo?
\choice
{$4$}
{$6$}
{$3$}
{\True $5$}
\loigiai{
Tổng số kẹo Nam đã ăn và cho em là: $2 + 3 = 5$ (chiếc kẹo).\nSố kẹo còn lại trong hộp là: $10 - 5 = 5$ (chiếc kẹo).
\textbf{Đáp số:} $5$.\par
\textbf{Đáp án của câu hỏi là D}
}
\end{ex}

\begin{ex} %[1D3V11]
Mẹ mua về 8 quả trứng. Bữa sáng mẹ dùng 3 quả, bữa tối mẹ dùng thêm 2 quả. Hỏi mẹ còn lại bao nhiêu quả trứng?
\choice
{\True $3$}
{$2$}
{$4$}
{$1$}
\loigiai{
Tổng số trứng mẹ đã dùng là: $3 + 2 = 5$ (quả).\nSố trứng còn lại là: $8 - 5 = 3$ (quả trứng).
\textbf{Đáp số:} $3$.\par
\textbf{Đáp án của câu hỏi là A}
}
\end{ex}

\begin{ex} %[1D3V11]
Bạn Hùng có 10 viên bi. Hùng cho bạn 4 viên bi, sau đó làm rơi mất 1 viên bi. Hỏi Hùng còn lại bao nhiêu viên bi?
\choice
{$4$}
{$6$}
{\True $5$}
{$3$}
\loigiai{
Số viên bi Hùng còn lại là: $10 - 4 - 1 = 5$ (viên bi).
\textbf{Đáp số:} $5$.\par
\textbf{Đáp án của câu hỏi là C}
}
\end{ex}

\begin{ex} %[1D3V11]
Tìm một số tự nhiên biết rằng lấy số lớn nhất có một chữ số trừ đi số đó thì được 6.
\choice
{$2$}
{\True $3$}
{$4$}
{$1$}
\loigiai{
Số lớn nhất có một chữ số là 9.\nSố cần tìm là: $9 - 6 = 3$.
\textbf{Đáp số:} $3$.\par
\textbf{Đáp án của câu hỏi là B}
}
\end{ex}

\begin{ex} %[1D3V11]
Tìm số tự nhiên $x$ thỏa mãn: $9 - x = 10 - 8$.
\choice
{$6$}
{$8$}
{$5$}
{\True $7$}
\loigiai{
Vế phải: $10 - 8 = 2$.\nTa có: $9 - x = 2 \Rightarrow x = 9 - 2 = 7$.
\textbf{Đáp số:} $7$.\par
\textbf{Đáp án của câu hỏi là D}
}
\end{ex}

\begin{ex} %[1D3V11]
Tìm số tự nhiên $y$ thỏa mãn: $y - 3 = 8 - 4$.
\choice
{\True $7$}
{$6$}
{$8$}
{$5$}
\loigiai{
Vế phải: $8 - 4 = 4$.\nTa có: $y - 3 = 4 \Rightarrow y = 4 + 3 = 7$.
\textbf{Đáp số:} $7$.\par
\textbf{Đáp án của câu hỏi là A}
}
\end{ex}

\textbf{Phần 2. Trắc nghiệm đúng sai}
\setcounterpageref{ex}{0}

\begin{ex} %[1D3V11]
Xét tính đúng hoặc sai của các khẳng định toán học sau đây:
\choiceTF
{\True Một sợi dây dài 10 cm. Bạn Mai cắt đi một đoạn dài 4 cm. đoạn dây còn lại dài $6\text{ cm}$ xăng-ti-mét.}
{Một cuốn sách có 10 trang. Bạn An đã đọc được 7 trang. An còn phải đọc $2$ trang nữa để hết cuốn sách.}
{Có 8 bạn nhỏ chơi trốn tìm. Đã tìm thấy 5 bạn. còn $2$ bạn chưa bị tìm thấy.}
{\True Kết quả của bài toán «Tìm hiệu giữa số bé nhất có hai chữ số và số lớn nhất có một chữ số bé hơn 7» là $4$.}
\loigiai{
\begin{itemize}
    \item Ý a) Đúng: Độ dài đoạn dây còn lại là: $10 - 4 = 6$ (cm).
\textbf{Đáp số:} $6$.
    \item Ý b) Sai: Kết quả đúng là $3$, không phải $2$. Số trang sách An còn phải đọc là: $10 - 7 = 3$ (trang).
\textbf{Đáp số:} $3$.
    \item Ý c) Sai: Kết quả đúng là $3$, không phải $2$. Số bạn chưa bị tìm thấy là: $8 - 5 = 3$ (bạn).
\textbf{Đáp số:} $3$.
    \item Ý d) Đúng: Số bé nhất có hai chữ số là $10$.\nSố lớn nhất có một chữ số bé hơn 7 là số $6$.\nHiệu của hai số là: $10 - 6 = 4$.
\textbf{Đáp số:} $4$.
\end{itemize}
\begin{itemize}
    \item a) \textbf{Đúng}.
    \item b) \textbf{Sai}.
    \item c) \textbf{Sai}.
    \item d) \textbf{Đúng}.
\end{itemize}
}
\end{ex}

\begin{ex} %[1D3V11]
Xét tính đúng hoặc sai của các bài toán và phép tính sau đây:
\choiceTF
{Kết quả của bài toán «Tìm một số biết rằng lấy 10 trừ đi số đó rồi trừ tiếp đi 2 thì được kết quả bằng 5» là $2$.}
{\True Bạn Lan có 9 cái kẹo, bạn Mai có ít hơn bạn Lan 4 cái kẹo. bạn Mai có $5$ cái kẹo.}
{\True Trong giỏ có 7 quả cam gồm 4 quả chín và số còn lại là quả xanh. số cam xanh ít hơn số cam chín $1$ quả.}
{Có 8 bạn nhỏ chơi trốn tìm. Đã tìm thấy 5 bạn. còn $2$ bạn chưa bị tìm thấy.}
\loigiai{
\begin{itemize}
    \item Ý a) Sai: Kết quả đúng là $3$, không phải $2$. Ta có: $10 - x - 2 = 5 \Rightarrow 10 - x = 7 \Rightarrow x = 10 - 7 = 3$.
\textbf{Đáp số:} $3$.
    \item Ý b) Đúng: Số cái kẹo của bạn Mai là: $9 - 4 = 5$ (cái kẹo).
\textbf{Đáp số:} $5$.
    \item Ý c) Đúng: Số quả cam xanh là: $7 - 4 = 3$ (quả).\nSố cam xanh ít hơn số cam chín là: $4 - 3 = 1$ (quả).
\textbf{Đáp số:} $1$.
    \item Ý d) Sai: Kết quả đúng là $3$, không phải $2$. Số bạn chưa bị tìm thấy là: $8 - 5 = 3$ (bạn).
\textbf{Đáp số:} $3$.
\end{itemize}
\begin{itemize}
    \item a) \textbf{Sai}.
    \item b) \textbf{Đúng}.
    \item c) \textbf{Đúng}.
    \item d) \textbf{Sai}.
\end{itemize}
}
\end{ex}

\textbf{Phần 3. Trả lời ngắn}
\setcounterpageref{ex}{0}

\begin{ex} %[1D3V11]
Một đàn vịt có 10 con. Có 4 con lên bờ nghỉ ngơi, 2 con khác bơi đi chỗ khác. Hỏi còn bao nhiêu con vịt bơi ở đó?
\shortans{4}
\loigiai{
\textbf{Đáp số:} $4$\par
Số con vịt còn bơi ở đó là: $10 - 4 - 2 = 4$ (con vịt).
\textbf{Đáp số:} $4$.
}
\end{ex}

\begin{ex} %[1D3V11]
Tìm số tự nhiên $m$ sao cho: $7 - m > 5$. Hỏi $m$ có thể nhận bao nhiêu giá trị tự nhiên?
\shortans{2}
\loigiai{
\textbf{Đáp số:} $2$\par
Các số tự nhiên $m$ sao cho $7 - m > 5$ là:\n- $m = 0$ (vì $7 - 0 = 7 > 5$)\n- $m = 1$ (vì $7 - 1 = 6 > 5$)\nVậy có $2$ giá trị thỏa mãn.
\textbf{Đáp số:} $2$.
}
\end{ex}

\begin{ex} %[1D3V11]
Một cây nến cao 9 cm, bị cháy ngắn đi 5 cm. Hỏi cây nến còn lại cao bao nhiêu xăng-ti-mét?
\shortans{4}
\loigiai{
\textbf{Đáp số:} $4$\par
Chiều cao còn lại của cây nến là: $9 - 5 = 4$ (cm).
\textbf{Đáp số:} $4$.
}
\end{ex}

\begin{ex} %[1D3V11]
Lấy 10 trừ đi số chẵn lớn nhất có một chữ số. Kết quả bằng bao nhiêu?
\shortans{2}
\loigiai{
\textbf{Đáp số:} $2$\par
Số chẵn lớn nhất có một chữ số là số $8$.\nLấy 10 trừ đi 8: $10 - 8 = 2$.
\textbf{Đáp số:} $2$.
}
\end{ex}

\begin{ex} %[1D3V11]
Lấy 9 trừ đi số lẻ bé nhất khác 0. Kết quả bằng bao nhiêu?
\shortans{8}
\loigiai{
\textbf{Đáp số:} $8$\par
Số lẻ bé nhất khác 0 là số $1$.\nLấy 9 trừ đi 1: $9 - 1 = 8$.
\textbf{Đáp số:} $8$.
}
\end{ex}

\begin{ex} %[1D3V11]
Bác An nuôi 10 con bồ câu. Buổi sáng bác mở chuồng cho 6 con bay ra ngoài kiếm ăn. Hỏi trong chuồng còn lại bao nhiêu con bồ câu?
\shortans{4}
\loigiai{
\textbf{Đáp số:} $4$\par
Số bồ câu còn lại trong chuồng là: $10 - 6 = 4$ (con bồ câu).
\textbf{Đáp số:} $4$.
}
\end{ex}
